\chapter{Python programs for exact arithmetic}
\label{appendix-software-resources}
\chapteroverview{Exact arithmetic separates algebraic identities from rounding
error. These programs construct finite fields, carry out modular Fourier
transforms, and implement the encoding and decoding procedures of
Chapter~\ref{chap-extension-finite-field-values}. Their small examples are
checked against direct transforms, integer products, and code constraints.}

The modules \texttt{finite\_fields.py} and
\texttt{error\_correcting\_codes.py} use SymPy for integer and polynomial
arithmetic, and \texttt{galois} for finite-field arrays. Their imports,
shared helpers, and executable checks are supplied with the complete Python
sources. Array coordinates increase with polynomial degree; a polynomial
library may instead print its coefficients in descending order, so the
conversion is explicit. All random examples use a fixed seed.

\section{Fourier transforms over finite fields}
\label{sect1-listing-transform-finite-body}
Program~\ref{listing-fft-finite-body} factors $\Phi_n$ over $\FF_p$ and
uses one irreducible factor to construct a field containing a root of
\emph{exact} order $n$. It rejects $p\mid n$. The example takes $p=2$
and $n=15$, so the transform values lie in an extension of the binary
input alphabet. The module also supplies a mixed-radix FFT, checked against
the direct transform here.

In \texttt{galois}, an integer used to construct an extension-field element
encodes its polynomial-basis coordinates. Consequently, the scalar $n\,1_K$
is constructed as \texttt{field(n \% p)}, not \texttt{field(n)}. This
matters in the inverse transform's normalization.
\begin{listing}[H]
\PythonCode{listing-fft-finite-body}
\caption{Constructing and checking an exact finite-field Fourier transform.}
\label{listing-fft-finite-body}
\end{listing}

\section{Fourier transforms over residue rings}
\label{sect1-listing-transform-ring}
Program~\ref{listing-fft-ring} uses Python integers reduced modulo $m$.
For a dyadic transform length $n>1$, it checks that $n$ is invertible and
$\zeta^{n/2}=-1\pmod m$. These hypotheses make $\zeta$ a principal root
by Proposition~\ref{prop-cns-square-root-ring}. A modulus need not be prime:
the transform works, for instance, modulo the composite Fermat number
$2^{32}+1$. The inverse uses the inverse root and multiplication by
$n^{-1}\pmod m$.
\begin{listing}[H]
\PythonCode{listing-fft-ring}
\caption{A radix-two Fourier transform with modular integer arithmetic.}
\label{listing-fft-ring}
\end{listing}

\section{Multiplication of large integers}
\label{sect1-listing-mult-large-integers-ring}
The digit base is an input parameter, and the transform length covers the
full linear convolution. Program~\ref{listing-mult-large-integers} chooses
a modulus $m=2^{cn/2}+1$ larger than the coefficient bound; $2^c$ is then
a principal $n$th root. The inverse transform returns the exact nonnegative
convolution coefficients, after which carry propagation recovers the product.
Signed inputs and zero are handled separately. The digit-conversion helpers
are given in Programs~\ref{listing-number2vector} and~\ref{listing-vector2number}.

Python's built-in integers already support arbitrary precision. The purpose
of this implementation is to expose the transform-based algorithm and its
bounds, rather than to extend Python's range or promise a faster product.
The tests compare the result with independent built-in multiplication,
including operands with hundreds of decimal digits.
\begin{listing}[H]
\PythonCode{listing-mult-large-integers}
\caption{Exact integer multiplication with a sufficiently large modular transform.}
\label{listing-mult-large-integers}
\end{listing}

\section{BCH encoding and decoding}
\label{sect1-listing-decoding-bch}
The implementation has four parts:
\begin{enumerate}
\item Program~\ref{listing-decodage-bch-1} takes the union of the cyclotomic
cosets meeting $1,\ldots,\delta-1$ and constructs the generator polynomial.
The alphabet is $\FF_p$, while root and syndrome computations take place
in a suitable extension.
\item Program~\ref{listing-decodage-bch-2} multiplies a message polynomial
by the generator, then illustrates correction of three errors in a binary
length-31 code of designed distance seven.
\item Program~\ref{listing-decodage-bch-3} determines the actual error
count from a Hankel matrix and solves for the minimal locator. This
includes the zero-error case and fewer errors than the permitted maximum.
\item Program~\ref{listing-decodage-bch-4} either completes the syndrome
spectrum or finds the locator roots and solves for error magnitudes.
Both routes verify the alphabet, error count, and all code constraints.
\end{enumerate}
The rank calculation is exact finite-field elimination. Outside the stated
correction radius, successful checks do not identify the transmitted word
uniquely; this is the usual limitation of bounded-distance decoding.
\begin{listing}[H]
\PythonCode{listing-decodage-bch-1}
\caption{Constructing the generator of a prime-alphabet BCH code.}
\label{listing-decodage-bch-1}
\end{listing}
\begin{listing}[H]
\PythonCode{listing-decodage-bch-2}
\caption{BCH encoding and an independently checked decoding example.}
\label{listing-decodage-bch-2}
\end{listing}
\begin{listing}[H]
\PythonCode{listing-decodage-bch-3}
\caption{Recovering the actual locator degree and its coefficients.}
\label{listing-decodage-bch-3}
\end{listing}
\begin{listing}[H]
\PythonCode{listing-decodage-bch-4}
\caption{BCH decoding by spectrum completion or locator roots.}
\label{listing-decodage-bch-4}
\end{listing}
